\subsection{Input buffer}
\begin{figure*}[h]\centering % Using \begin{figure*} makes the figure take up the entire width of the page
	\includegraphics[width=\linewidth]{figure3}
	\caption{Input buffer}
	\label{fig:figure3}
\end{figure*}

The block diagram of the input buffer module is shown in Figure \ref{fig:figure3}. BRAM 0 is
used to store complete even frames and BRAM 1 stores complete odd frames, where
a frame is an array of pixels with shape [24,32,1]. All values are 18 bit
signed integers. These BRAMs contain 2 write ports and 2 read ports. One write
port is used by the I2C FSM to store and deinterlace incoming pixels from the
sensor. One read port is used by the AXI-S FSM to read pixels in raster order
and send them out the AXI-S interface. The second read and write ports are used
to provide a memory mapped interface to the pixel values.
The input buffer module contains the following memory map which is accessible through the AXI memory mapped slave interface.
\begin{table}[hbt]
\caption{Input buffer memory map}
%	\centering
\begin{tabular}{|l|r|l|l|l|}
\toprule
%		\multicolumn{2}{c}{Name} \\
%		\cmidrule(r){1-2}
Address & Bits & RW & Reset & Description \\
\midrule
	\hline
	0x0	&0	&W	&1	&Bit bang SDA out (open drain) \\
	\hline
	0x0	&1	&R	&-	&Bit bang SDA in \\
	\hline
	0x0	&2	&W	&1	&Bit bang SCL out (open drain) \\
	\hline
	0x0	&3	&R	&-	&Bit bang SCL in \\
	\hline
	0x4	&0	&R/W	&0	&Run/Stop \\
	\hline
	0x1000	&17:0	&R/W		&&BRAM 2 EEPROM calibration coefficients for pixel compensation \\
	\hline
	0x2000	&17:0	&R/W		&&BRAM 0 SoC CPU access \\
	\hline
	0x3000	&17:0	&R/W		&&BRAM 1 SoC CPU access \\
\bottomrule
\end{tabular}
\label{tab:mmap}
\end{table}

\paragraph{I2C FSM}
The ibuf.v module contains an I2C master FSM which periodically reads IR pixel
values from the sensor, applies per-pixel compensation factors, and stores the
pixels in an on-chip 32x24 BRAM frame buffer. Refer to the
\href{https://www.melexis.com/en/documents/documentation/datasheets/datasheet-mlx90640}{MLX90640 data sheet}.
for details on sensor initialization, reading pixel values, and applying
compensation factors.  The bit bang interface is used by the SoC application
CPU to initialize the MLX90640 sensor. The sensor should be initialized with
factory default settings except for the IR refresh rate which should be set to
64Hz. Note that 64Hz is the interlaced rate for one subframe, the actual frame
rate is 32Hz. It is desirable to make the refresh rate configurable in order to
trade off thermal noise against refresh rate. At initialization, the
calibration parameters must be read from the MLX90640 EEPROM using I2C reads,
mathematically combined into 768 per-pixel coefficients, and stored in BRAM 2.
Initially these coefficients can be set to unity to bypass calibration
compensation. The compensation procedure is described in section 11.2.2 of the
MLX90640 specification. Pixel compensation will require the use of a DSP
multiplier and fixed point arithmetic.
For simplicity, the input buffer can be double buffered so that one BRAM is
receiving from the I2C FSM while the other is sending to the AXI-S FSM. Note
that a 768-pixel frame fits in one 18x1K BRAM. The overall pseudo code for the
I2C FSM follows.
\begin{enumerate}[noitemsep] % [noitemsep] removes whitespace between the items for a compact look
\item	Poll MLX90640 until new data available
\item	Read subframe 0 pixels from sensor RAM starting at 0x400
\item	Multiply each pixel by its individual compensation coefficient (768 fixed point coefficients stored in BRAM 2)
\item	Store compensated pixels in BRAM 0 (even frames)
\item	Repeat 1-4 for subframe 1 (chess pattern interleave)
\item	Assert interrupt to SoC CPU, trigger AXI-S FSM
\item	Repeat 1-6 for BRAM 1 (odd frames)
\end{enumerate}

Note that if the input pixels arrive in raster order there is no need for a
frame buffer. However, the MLX90640 sensor captures pixels in an interlaced
chess board format. The frame buffer is designed to deinterlace the two chess
boards into a single complete frame and then sequentially send each pixel over
the output AXI-S interface in raster order. The frame buffer uses double
buffered BRAM so the I2C side and AXI-S side can operate simultaneously.
Additionally, the ibuf.v module generates an interrupt to the SoC CPU when a
complete frame arrives (at 32Hz), so that the CPU can read the 768 raw pixel
values in real time.

\paragraph{AXI-S FSM}
The AXI-S FSM operates concurrently with the I2C FSM and is responsible for
emitting frames over the AXI stream interface. The overall pseudo code for the
AXI-S FSM follows.
\begin{enumerate}[noitemsep] % [noitemsep] removes whitespace between the items for a compact look
	\item	Receive frame 0 trigger signal from I2C FSM \verb|foo_bar|
\item	Sequentially read pixels from BRAM 0 in raster order and output to AXI-S
\item	Rate limit AXI-S so that TVALID toggles at ~25KHz (768 pixels * 32Hz)
\item	Assert TLAST
\item	Repeat 1-4 for frame 1 / BRAM 1
\end{enumerate}
Conceptually, data on the AXI-S interface is a tensor sent sequentially in 'C'
order and framed using TLAST. The shape of the IR sensor data is [24,32,1]. The
pixels should be in raster order with 0,0 in the upper left corner of the
frame.

